\chapter{附录}
\label{ch:appendix}

% =============================================================
% Appendix A: Complete Company Index
% Appendix B: Glossary of Terms
% Appendix C: Data Sources and Methodology
% Appendix D: Cross-Reference to Emergence
% =============================================================


% =============================================================
%  §A.1  COMPLETE COMPANY INDEX
% =============================================================
\section{完整公司索引}
\label{sec:appendix-index}

以下按拉丁字母顺序列出本书提及的主要 AI 创业公司及相关企业。
每条记录包含公司名（英文及中文，如适用）、总部所在地、成立年份、
所属赛道/层级，以及本书主要讨论章节。
表中"层级"一栏的缩写含义如下：\textbf{基}=基础设施层，
\textbf{模}=模型层，\textbf{平}=平台与中间件层，
\textbf{工}=开发者工具层，\textbf{C}=C~端应用层，
\textbf{B}=B~端/企业应用层，\textbf{跨}=跨层级。

\medskip

\noindent
\textbf{说明：}成立年份以公司实体注册或首次公开运营的年份为准；
总部所在地取截至2026年3月的信息。
部分公司因业务跨层级，在多个章节中出现。

\bigskip

% --- A ---
\noindent{\large\bfseries A}\nopagebreak
\begin{tabularx}{\textwidth}{@{}l X l l l l@{}}
\toprule
\rowcolor{figLightGray}
\textbf{公司名} & \textbf{中文名} & \textbf{总部} & \textbf{成立} & \textbf{层级} & \textbf{章节} \\
\midrule
Ada Health          & ---             & Berlin        & 2011 & C    & \ref{ch:consumer-life} \\
AI21 Labs           & ---             & Tel Aviv      & 2017 & 模   & \ref{ch:models}, \ref{ch:global} \\
Anyscale            & ---             & San Francisco & 2019 & 基   & \ref{ch:infrastructure} \\
Anthropic           & ---             & San Francisco & 2021 & 模   & \ref{ch:models}, \ref{ch:funding} \\
Astera Labs         & ---             & Santa Clara   & 2017 & 基   & \ref{ch:infrastructure} \\
Augment Code        & ---             & Palo Alto     & 2022 & 工   & \ref{ch:devtools} \\
Ayar Labs           & ---             & San Jose      & 2015 & 基   & \ref{ch:infrastructure} \\
\bottomrule
\end{tabularx}

\bigskip

% --- B ---
\noindent{\large\bfseries B}\nopagebreak
\begin{tabularx}{\textwidth}{@{}l X l l l l@{}}
\toprule
\rowcolor{figLightGray}
\textbf{公司名} & \textbf{中文名} & \textbf{总部} & \textbf{成立} & \textbf{层级} & \textbf{章节} \\
\midrule
Baichuan (百川智能) & 百川智能       & Beijing       & 2023 & 模   & \ref{ch:models}, \ref{ch:china} \\
Baseten             & ---            & San Francisco & 2019 & 基   & \ref{ch:infrastructure} \\
Buildots            & ---            & Tel Aviv      & 2018 & B    & \ref{ch:enterprise} \\
\bottomrule
\end{tabularx}

\bigskip

% --- C ---
\noindent{\large\bfseries C}\nopagebreak
\begin{tabularx}{\textwidth}{@{}l X l l l l@{}}
\toprule
\rowcolor{figLightGray}
\textbf{公司名} & \textbf{中文名} & \textbf{总部} & \textbf{成立} & \textbf{层级} & \textbf{章节} \\
\midrule
Cerebras Systems    & ---             & Sunnyvale     & 2016 & 基   & \ref{ch:infrastructure} \\
Character.ai        & ---             & Menlo Park    & 2021 & C    & \ref{ch:consumer-social} \\
Chroma              & ---             & San Francisco & 2022 & 平   & \ref{ch:platform} \\
Claude Code (Anthropic) & ---         & San Francisco & 2025 & 工   & \ref{ch:devtools} \\
Clay                & ---             & New York      & 2023 & B    & \ref{ch:consumer-social} \\
Cleo                & ---             & London        & 2016 & C    & \ref{ch:consumer-life} \\
Cognition (Devin)   & ---             & San Francisco & 2023 & 工   & \ref{ch:devtools} \\
Cohere              & ---             & Toronto       & 2019 & 模   & \ref{ch:models}, \ref{ch:global} \\
CoreWeave           & ---             & Roseland, NJ  & 2017 & 基   & \ref{ch:infrastructure}, \ref{ch:funding} \\
CrowdStrike         & ---             & Austin        & 2011 & B    & \ref{ch:enterprise} \\
Crusoe Energy       & ---             & San Francisco & 2020 & 基   & \ref{ch:infrastructure} \\
Cursor (Anysphere)  & ---             & San Francisco & 2022 & 工   & \ref{ch:devtools}, \ref{ch:funding} \\
\bottomrule
\end{tabularx}

\bigskip

% --- D ---
\noindent{\large\bfseries D}\nopagebreak
\begin{tabularx}{\textwidth}{@{}l X l l l l@{}}
\toprule
\rowcolor{figLightGray}
\textbf{公司名} & \textbf{中文名} & \textbf{总部} & \textbf{成立} & \textbf{层级} & \textbf{章节} \\
\midrule
d-Matrix            & ---             & Santa Clara   & 2019 & 基   & \ref{ch:infrastructure} \\
DeepL               & ---             & Cologne       & 2017 & C    & \ref{ch:consumer-prod} \\
DeepSeek (深度求索) & 深度求索        & Hangzhou      & 2023 & 模   & \ref{ch:models}, \ref{ch:china} \\
Dify                & ---             & San Francisco & 2023 & 平   & \ref{ch:platform} \\
Duolingo            & ---             & Pittsburgh    & 2011 & C    & \ref{ch:consumer-life} \\
\bottomrule
\end{tabularx}

\bigskip

% --- E ---
\noindent{\large\bfseries E}\nopagebreak
\begin{tabularx}{\textwidth}{@{}l X l l l l@{}}
\toprule
\rowcolor{figLightGray}
\textbf{公司名} & \textbf{中文名} & \textbf{总部} & \textbf{成立} & \textbf{层级} & \textbf{章节} \\
\midrule
Eightfold AI        & ---             & Santa Clara   & 2016 & B    & \ref{ch:consumer-social} \\
ElevenLabs          & ---             & New York      & 2022 & 模/C & \ref{ch:consumer-create} \\
Elicit              & ---             & San Francisco & 2021 & C    & \ref{ch:consumer-prod} \\
Enfabrica           & ---             & San Jose      & 2020 & 基   & \ref{ch:infrastructure} \\
Etched              & ---             & Cupertino     & 2022 & 基   & \ref{ch:infrastructure} \\
EvenUp              & ---             & San Francisco & 2019 & C    & \ref{ch:consumer-life} \\
\bottomrule
\end{tabularx}

\bigskip

% --- F ---
\noindent{\large\bfseries F}\nopagebreak
\begin{tabularx}{\textwidth}{@{}l X l l l l@{}}
\toprule
\rowcolor{figLightGray}
\textbf{公司名} & \textbf{中文名} & \textbf{总部} & \textbf{成立} & \textbf{层级} & \textbf{章节} \\
\midrule
Fireworks AI        & ---             & San Francisco & 2022 & 基/平 & \ref{ch:infrastructure} \\
FluidStack          & ---             & London        & 2020 & 基    & \ref{ch:infrastructure} \\
Freeletics          & ---             & Munich        & 2013 & C     & \ref{ch:consumer-life} \\
\bottomrule
\end{tabularx}

\bigskip

% --- G ---
\noindent{\large\bfseries G}\nopagebreak
\begin{tabularx}{\textwidth}{@{}l X l l l l@{}}
\toprule
\rowcolor{figLightGray}
\textbf{公司名} & \textbf{中文名} & \textbf{总部} & \textbf{成立} & \textbf{层级} & \textbf{章节} \\
\midrule
G42                 & ---             & Abu Dhabi     & 2018 & 跨   & \ref{ch:global} \\
Gamma               & ---             & San Francisco & 2020 & C    & \ref{ch:consumer-prod} \\
GitHub Copilot (Microsoft) & ---     & San Francisco & 2021 & 工   & \ref{ch:devtools} \\
Glean               & ---             & Palo Alto     & 2019 & B    & \ref{ch:enterprise} \\
Gong                & ---             & San Francisco & 2015 & B    & \ref{ch:consumer-social} \\
Google DeepMind     & ---             & London        & 2010 & 模   & \ref{ch:models} \\
Groq                & ---             & Mountain View & 2016 & 基   & \ref{ch:infrastructure} \\
\bottomrule
\end{tabularx}

\bigskip

% --- H ---
\noindent{\large\bfseries H}\nopagebreak
\begin{tabularx}{\textwidth}{@{}l X l l l l@{}}
\toprule
\rowcolor{figLightGray}
\textbf{公司名} & \textbf{中文名} & \textbf{总部} & \textbf{成立} & \textbf{层级} & \textbf{章节} \\
\midrule
Hailo               & ---             & Tel Aviv      & 2017 & 基   & \ref{ch:infrastructure}, \ref{ch:global} \\
Harvey              & ---             & San Francisco & 2022 & 模/C & \ref{ch:models}, \ref{ch:consumer-life} \\
HeyGen              & ---             & Los Angeles   & 2020 & 模/C & \ref{ch:consumer-create} \\
Hippocratic AI      & ---             & Palo Alto     & 2023 & 模   & \ref{ch:models} \\
Horizon Robotics (地平线) & 地平线   & Beijing       & 2015 & 基   & \ref{ch:infrastructure}, \ref{ch:china} \\
Hugging Face        & ---             & New York      & 2016 & 平   & \ref{ch:platform} \\
\bottomrule
\end{tabularx}

\bigskip

% --- J ---
\noindent{\large\bfseries J}\nopagebreak
\begin{tabularx}{\textwidth}{@{}l X l l l l@{}}
\toprule
\rowcolor{figLightGray}
\textbf{公司名} & \textbf{中文名} & \textbf{总部} & \textbf{成立} & \textbf{层级} & \textbf{章节} \\
\midrule
Jiepu Xingchen (阶跃星辰) & 阶跃星辰 & Beijing    & 2023 & 模   & \ref{ch:models}, \ref{ch:china} \\
\bottomrule
\end{tabularx}

\bigskip

% --- K ---
\noindent{\large\bfseries K}\nopagebreak
\begin{tabularx}{\textwidth}{@{}l X l l l l@{}}
\toprule
\rowcolor{figLightGray}
\textbf{公司名} & \textbf{中文名} & \textbf{总部} & \textbf{成立} & \textbf{层级} & \textbf{章节} \\
\midrule
K Health            & ---             & New York      & 2016 & C    & \ref{ch:consumer-life} \\
Khan Academy        & ---             & Mountain View & 2008 & C    & \ref{ch:consumer-life} \\
Kling AI (可灵)     & 可灵            & Beijing       & 2024 & C    & \ref{ch:consumer-create}, \ref{ch:china} \\
Kunlunxin (昆仑芯)  & 昆仑芯          & Beijing       & 2021 & 基   & \ref{ch:infrastructure}, \ref{ch:china} \\
\bottomrule
\end{tabularx}

\bigskip

% --- L ---
\noindent{\large\bfseries L}\nopagebreak
\begin{tabularx}{\textwidth}{@{}l X l l l l@{}}
\toprule
\rowcolor{figLightGray}
\textbf{公司名} & \textbf{中文名} & \textbf{总部} & \textbf{成立} & \textbf{层级} & \textbf{章节} \\
\midrule
Lakera              & ---             & Zurich        & 2023 & 平   & \ref{ch:platform} \\
Lambda              & ---             & San Francisco & 2012 & 基   & \ref{ch:infrastructure} \\
LangChain           & ---             & San Francisco & 2022 & 平   & \ref{ch:platform} \\
Langfuse            & ---             & Berlin        & 2023 & 平   & \ref{ch:platform} \\
Lightmatter         & ---             & Boston        & 2017 & 基   & \ref{ch:infrastructure} \\
Lightning AI        & ---             & New York      & 2019 & 基   & \ref{ch:infrastructure} \\
LlamaIndex          & ---             & San Francisco & 2022 & 平   & \ref{ch:platform} \\
Luma AI             & ---             & San Francisco & 2021 & 模/C & \ref{ch:consumer-create} \\
\bottomrule
\end{tabularx}

\bigskip

% --- M ---
\noindent{\large\bfseries M}\nopagebreak
\begin{tabularx}{\textwidth}{@{}l X l l l l@{}}
\toprule
\rowcolor{figLightGray}
\textbf{公司名} & \textbf{中文名} & \textbf{总部} & \textbf{成立} & \textbf{层级} & \textbf{章节} \\
\midrule
MatX                & ---             & Mountain View & 2022 & 基   & \ref{ch:infrastructure} \\
Meta AI             & ---             & Menlo Park    & 2013 & 模   & \ref{ch:models} \\
Microsoft Copilot   & ---             & Redmond       & 2023 & B    & \ref{ch:enterprise} \\
Midjourney          & ---             & San Francisco & 2022 & C    & \ref{ch:consumer-create} \\
Milvus / Zilliz     & ---             & Redwood City  & 2017 & 平   & \ref{ch:platform} \\
MiniMax (稀宇科技)  & 稀宇科技        & Shanghai      & 2021 & 模   & \ref{ch:models}, \ref{ch:china} \\
Mintlify            & ---             & San Francisco & 2022 & 工   & \ref{ch:devtools} \\
Mistral             & ---             & Paris         & 2023 & 模   & \ref{ch:models}, \ref{ch:global} \\
Modal               & ---             & San Francisco & 2021 & 基   & \ref{ch:infrastructure} \\
Momentic            & ---             & San Francisco & 2023 & 工   & \ref{ch:devtools} \\
Monarch Money       & ---             & San Francisco & 2019 & C    & \ref{ch:consumer-life} \\
Moonshot AI (月之暗面) & 月之暗面     & Beijing       & 2023 & 模   & \ref{ch:models}, \ref{ch:china} \\
Moore Threads (摩尔线程) & 摩尔线程   & Beijing       & 2020 & 基   & \ref{ch:infrastructure}, \ref{ch:china} \\
Muxi (沐曦)         & 沐曦            & Shanghai      & 2020 & 基   & \ref{ch:infrastructure}, \ref{ch:china} \\
\bottomrule
\end{tabularx}

\bigskip

% --- N ---
\noindent{\large\bfseries N}\nopagebreak
\begin{tabularx}{\textwidth}{@{}l X l l l l@{}}
\toprule
\rowcolor{figLightGray}
\textbf{公司名} & \textbf{中文名} & \textbf{总部} & \textbf{成立} & \textbf{层级} & \textbf{章节} \\
\midrule
Nebius              & ---             & Amsterdam     & 2024 & 基   & \ref{ch:infrastructure} \\
Notion AI           & ---             & San Francisco & 2013 & C    & \ref{ch:consumer-prod} \\
\bottomrule
\end{tabularx}

\bigskip

% --- O ---
\noindent{\large\bfseries O}\nopagebreak
\begin{tabularx}{\textwidth}{@{}l X l l l l@{}}
\toprule
\rowcolor{figLightGray}
\textbf{公司名} & \textbf{中文名} & \textbf{总部} & \textbf{成立} & \textbf{层级} & \textbf{章节} \\
\midrule
Ollama              & ---             & San Francisco & 2023 & 平   & \ref{ch:platform} \\
OpenAI              & ---             & San Francisco & 2015 & 模   & \ref{ch:models}, \ref{ch:funding} \\
OpenRouter          & ---             & San Francisco & 2023 & 平   & \ref{ch:platform} \\
\bottomrule
\end{tabularx}

\bigskip

% --- P ---
\noindent{\large\bfseries P}\nopagebreak
\begin{tabularx}{\textwidth}{@{}l X l l l l@{}}
\toprule
\rowcolor{figLightGray}
\textbf{公司名} & \textbf{中文名} & \textbf{总部} & \textbf{成立} & \textbf{层级} & \textbf{章节} \\
\midrule
Perplexity          & ---             & San Francisco & 2022 & C    & \ref{ch:consumer-social} \\
Pika                & ---             & Palo Alto     & 2023 & 模/C & \ref{ch:consumer-create} \\
Pinecone            & ---             & San Francisco & 2019 & 平   & \ref{ch:platform} \\
\bottomrule
\end{tabularx}

\bigskip

% --- Q ---
\noindent{\large\bfseries Q}\nopagebreak
\begin{tabularx}{\textwidth}{@{}l X l l l l@{}}
\toprule
\rowcolor{figLightGray}
\textbf{公司名} & \textbf{中文名} & \textbf{总部} & \textbf{成立} & \textbf{层级} & \textbf{章节} \\
\midrule
QA Wolf             & ---             & San Francisco & 2019 & 工   & \ref{ch:devtools} \\
Qdrant              & ---             & Berlin        & 2021 & 平   & \ref{ch:platform} \\
\bottomrule
\end{tabularx}

\bigskip

% --- R ---
\noindent{\large\bfseries R}\nopagebreak
\begin{tabularx}{\textwidth}{@{}l X l l l l@{}}
\toprule
\rowcolor{figLightGray}
\textbf{公司名} & \textbf{中文名} & \textbf{总部} & \textbf{成立} & \textbf{层级} & \textbf{章节} \\
\midrule
Replika             & ---             & San Francisco & 2017 & C    & \ref{ch:consumer-social} \\
RunPod              & ---             & Cherry Hill, NJ & 2022 & 基 & \ref{ch:infrastructure} \\
Runway              & ---             & New York      & 2018 & 模/C & \ref{ch:consumer-create} \\
\bottomrule
\end{tabularx}

\bigskip

% --- S ---
\noindent{\large\bfseries S}\nopagebreak
\begin{tabularx}{\textwidth}{@{}l X l l l l@{}}
\toprule
\rowcolor{figLightGray}
\textbf{公司名} & \textbf{中文名} & \textbf{总部} & \textbf{成立} & \textbf{层级} & \textbf{章节} \\
\midrule
Salesforce          & ---             & San Francisco & 1999 & B    & \ref{ch:enterprise} \\
SambaNova Systems   & ---             & Palo Alto     & 2017 & 基   & \ref{ch:infrastructure} \\
Scale AI            & ---             & San Francisco & 2016 & 平   & \ref{ch:platform} \\
SF Compute          & ---             & San Francisco & 2023 & 基   & \ref{ch:infrastructure} \\
Sophgo (算能)       & 算能            & Beijing       & 2016 & 基   & \ref{ch:infrastructure}, \ref{ch:china} \\
Speak               & ---             & San Francisco & 2016 & C    & \ref{ch:consumer-life} \\
Stability AI        & ---             & London        & 2019 & 模   & \ref{ch:consumer-create} \\
Suno                & ---             & Cambridge, MA & 2023 & 模/C & \ref{ch:consumer-create} \\
Suiyuan (燧原科技)  & 燧原科技        & Shanghai      & 2018 & 基   & \ref{ch:infrastructure}, \ref{ch:china} \\
Synthesia           & ---             & London        & 2017 & 模/C & \ref{ch:consumer-create}, \ref{ch:global} \\
\bottomrule
\end{tabularx}

\bigskip

% --- T ---
\noindent{\large\bfseries T}\nopagebreak
\begin{tabularx}{\textwidth}{@{}l X l l l l@{}}
\toprule
\rowcolor{figLightGray}
\textbf{公司名} & \textbf{中文名} & \textbf{总部} & \textbf{成立} & \textbf{层级} & \textbf{章节} \\
\midrule
Tabnine             & ---             & Tel Aviv      & 2018 & 工   & \ref{ch:devtools}, \ref{ch:global} \\
Tenstorrent         & ---             & Toronto       & 2016 & 基   & \ref{ch:infrastructure} \\
Tianshu Zhixin (天数智芯) & 天数智芯  & Beijing       & 2018 & 基   & \ref{ch:infrastructure}, \ref{ch:china} \\
Together AI         & ---             & San Francisco & 2022 & 基/平 & \ref{ch:infrastructure} \\
Twelve Labs         & ---             & San Francisco & 2021 & 模   & \ref{ch:models}, \ref{ch:global} \\
\bottomrule
\end{tabularx}

\bigskip

% --- U ---
\noindent{\large\bfseries U}\nopagebreak
\begin{tabularx}{\textwidth}{@{}l X l l l l@{}}
\toprule
\rowcolor{figLightGray}
\textbf{公司名} & \textbf{中文名} & \textbf{总部} & \textbf{成立} & \textbf{层级} & \textbf{章节} \\
\midrule
UiPath              & ---             & New York      & 2005 & B    & \ref{ch:enterprise} \\
Upstage             & ---             & Seoul         & 2020 & 模   & \ref{ch:global} \\
\bottomrule
\end{tabularx}

\bigskip

% --- V ---
\noindent{\large\bfseries V}\nopagebreak
\begin{tabularx}{\textwidth}{@{}l X l l l l@{}}
\toprule
\rowcolor{figLightGray}
\textbf{公司名} & \textbf{中文名} & \textbf{总部} & \textbf{成立} & \textbf{层级} & \textbf{章节} \\
\midrule
Vast.ai             & ---             & San Francisco & 2020 & 基   & \ref{ch:infrastructure} \\
Vercel              & ---             & San Francisco & 2015 & 工   & \ref{ch:devtools} \\
\bottomrule
\end{tabularx}

\bigskip

% --- W ---
\noindent{\large\bfseries W}\nopagebreak
\begin{tabularx}{\textwidth}{@{}l X l l l l@{}}
\toprule
\rowcolor{figLightGray}
\textbf{公司名} & \textbf{中文名} & \textbf{总部} & \textbf{成立} & \textbf{层级} & \textbf{章节} \\
\midrule
Weaviate            & ---             & Amsterdam     & 2019 & 平   & \ref{ch:platform} \\
Windsurf / Codeium  & ---             & Mountain View & 2021 & 工   & \ref{ch:devtools} \\
Woebot Health       & ---             & San Francisco & 2017 & C    & \ref{ch:consumer-life} \\
Writer              & ---             & San Francisco & 2020 & 模   & \ref{ch:models} \\
Wysa                & ---             & Boston        & 2015 & C    & \ref{ch:consumer-life} \\
\bottomrule
\end{tabularx}

\bigskip

% --- X ---
\noindent{\large\bfseries X}\nopagebreak
\begin{tabularx}{\textwidth}{@{}l X l l l l@{}}
\toprule
\rowcolor{figLightGray}
\textbf{公司名} & \textbf{中文名} & \textbf{总部} & \textbf{成立} & \textbf{层级} & \textbf{章节} \\
\midrule
xAI                 & ---             & San Francisco & 2023 & 模   & \ref{ch:models}, \ref{ch:funding} \\
\bottomrule
\end{tabularx}

\bigskip

% --- Y ---
\noindent{\large\bfseries Y}\nopagebreak
\begin{tabularx}{\textwidth}{@{}l X l l l l@{}}
\toprule
\rowcolor{figLightGray}
\textbf{公司名} & \textbf{中文名} & \textbf{总部} & \textbf{成立} & \textbf{层级} & \textbf{章节} \\
\midrule
Yi / 01.AI (零一万物) & 零一万物      & Beijing       & 2023 & 模   & \ref{ch:models}, \ref{ch:china} \\
\bottomrule
\end{tabularx}

\bigskip

% --- Z ---
\noindent{\large\bfseries Z}\nopagebreak
\begin{tabularx}{\textwidth}{@{}l X l l l l@{}}
\toprule
\rowcolor{figLightGray}
\textbf{公司名} & \textbf{中文名} & \textbf{总部} & \textbf{成立} & \textbf{层级} & \textbf{章节} \\
\midrule
Zhipu AI (智谱 AI)  & 智谱 AI         & Beijing       & 2019 & 模   & \ref{ch:models}, \ref{ch:china} \\
\bottomrule
\end{tabularx}

\bigskip

% --- Chinese companies index (sorted by pinyin) ---
\noindent{\large\bfseries 中国 AI 芯片公司补充索引}\nopagebreak

\smallskip

以下公司在正文"中国 AI 芯片"部分（第~\ref{ch:infrastructure}~章§2.2b）
及第~\ref{ch:china}~章中有详细讨论。

\begin{tabularx}{\textwidth}{@{}l X l l l l@{}}
\toprule
\rowcolor{figLightGray}
\textbf{公司名} & \textbf{中文名} & \textbf{总部} & \textbf{成立} & \textbf{层级} & \textbf{章节} \\
\midrule
Biren Tech (壁仞科技)     & 壁仞科技   & Shanghai & 2019 & 基 & \ref{ch:infrastructure}, \ref{ch:china} \\
Cambricon (寒武纪)        & 寒武纪     & Beijing  & 2016 & 基 & \ref{ch:infrastructure}, \ref{ch:china} \\
Hygon (海光信息)          & 海光信息   & Chengdu  & 2014 & 基 & \ref{ch:infrastructure}, \ref{ch:china} \\
Kunlunxin (昆仑芯)        & 昆仑芯     & Beijing  & 2021 & 基 & \ref{ch:infrastructure}, \ref{ch:china} \\
Moore Threads (摩尔线程)   & 摩尔线程   & Beijing  & 2020 & 基 & \ref{ch:infrastructure}, \ref{ch:china} \\
Muxi (沐曦)               & 沐曦       & Shanghai & 2020 & 基 & \ref{ch:infrastructure}, \ref{ch:china} \\
Sophgo (算能)             & 算能       & Beijing  & 2016 & 基 & \ref{ch:infrastructure}, \ref{ch:china} \\
Suiyuan (燧原科技)        & 燧原科技   & Shanghai & 2018 & 基 & \ref{ch:infrastructure}, \ref{ch:china} \\
Tianshu Zhixin (天数智芯)  & 天数智芯   & Beijing  & 2018 & 基 & \ref{ch:infrastructure}, \ref{ch:china} \\
\bottomrule
\end{tabularx}

\medskip

\noindent\textit{注：本索引共收录约100家公司。因篇幅限制，部分在正文中
仅简要提及的公司（如各章节对比表中的次要案例）未逐一列出。
如需完整的公司--章节映射，请参考本书电子版的交互式索引。}


% =============================================================
%  §A.2  GLOSSARY OF TERMS
% =============================================================
\clearpage
\section{术语表}
\label{sec:appendix-glossary}

以下按英文字母顺序排列本书中出现的关键 AI/ML 术语，
提供中英对照名称及简明定义。

\bigskip

\begin{description}[style=nextline, leftmargin=2.5cm, labelwidth=2.3cm, font=\bfseries]

\item[AGI]
\textbf{通用人工智能}（Artificial General Intelligence）。
指具备人类水平通用推理能力的 AI 系统，能够在无需专门训练的情况下
完成各种智力任务。截至2026年，AGI 尚未实现，但各实验室已将其作为
公开目标。参见第~\ref{ch:outlook}~章讨论。

\item[Agent]
\textbf{AI 智能体/代理}。
一种能够自主规划、调用工具、执行多步任务的 AI 系统。
与简单的聊天机器人不同，Agent 具备目标导向的推理能力和
环境交互能力。2025--2026年的核心趋势之一。
参见第~\ref{ch:devtools}~章和第~\ref{ch:platform}~章。

\item[Alignment]
\textbf{对齐}。
确保 AI 系统的行为与人类意图和价值观一致的技术研究领域。
包括 RLHF、Constitutional AI、DPO 等方法。
Anthropic 和 OpenAI 的核心研究方向。

\item[API]
\textbf{应用程序编程接口}（Application Programming Interface）。
允许不同软件系统之间通信的标准化接口。在 AI 语境中，通常指
模型厂商提供的推理接口（如 OpenAI API），是整个 AI 应用层的
基础连接方式。

\item[ASIC]
\textbf{专用集成电路}（Application-Specific Integrated Circuit）。
为特定计算任务（如 AI 推理）专门设计的芯片，性能功耗比高于通用 GPU，
但灵活性较低。代表公司：Google TPU、Cerebras、Etched。
参见第~\ref{ch:infrastructure}~章。

\item[Attention]
\textbf{注意力机制}（Attention Mechanism）。
Transformer 架构的核心组件，允许模型在处理序列时动态关注不同位置的
信息。自注意力（Self-Attention）的计算复杂度为 $O(n^2)$，
是长上下文扩展的主要瓶颈。
参见 \textit{Emergence} 第2章。

\item[Benchmark]
\textbf{基准测试}。
用于评估模型能力的标准化测试集。常见基准包括 MMLU（通用知识）、
HumanEval（代码）、SWE-bench（软件工程）、MATH（数学推理）等。
参见第~\ref{ch:devtools}~章 SWE-bench 分析。

\item[Chain-of-Thought (CoT)]
\textbf{思维链}。
一种提示技术，引导模型在给出最终答案前展示中间推理步骤，
显著提升复杂推理任务的准确率。
是 o1/o3 等推理模型的基础技术。
参见 \textit{Emergence} 第10章。

\item[Constitutional AI]
\textbf{宪法 AI}。
Anthropic 提出的对齐方法，使用一组明确的原则（"宪法"）来指导
模型的自我改进，减少对人类反馈数据的依赖。
Claude 系列模型的核心技术之一。

\item[Context Window]
\textbf{上下文窗口}。
模型单次推理能处理的最大 token 数量。从 GPT-3 的 4K 到 2026 年
主流模型的 128K--1M+，上下文长度的扩展是过去三年最重要的
技术进步之一。参见 \textit{Emergence} 第7章。

\item[CUDA]
\textbf{统一计算设备架构}（Compute Unified Device Architecture）。
NVIDIA 的 GPU 并行计算平台和编程模型，经过十余年积累已成为深度学习
事实标准。CUDA 生态的锁定效应是 AI 芯片创业公司面临的最大壁垒。
参见第~\ref{ch:infrastructure}~章。

\item[Diffusion Model]
\textbf{扩散模型}。
一类通过学习"去噪"过程来生成数据的生成模型。
从纯噪声逐步还原出目标图像/视频/音频。
Stable Diffusion、DALL-E、Sora 的基础架构。
参见第~\ref{ch:consumer-create}~章和 \textit{Emergence} 相关讨论。

\item[Distillation]
\textbf{蒸馏}。
将大型"教师"模型的知识压缩到小型"学生"模型中的技术。
使小模型以更低成本接近大模型的性能。
DeepSeek-R1 蒸馏系列是2025年最成功的案例之一。

\item[DPO]
\textbf{直接偏好优化}（Direct Preference Optimization）。
一种替代 RLHF 的对齐方法，直接从人类偏好数据优化模型策略，
无需训练单独的奖励模型。算法更简单、训练更稳定。
参见 \textit{Emergence} 第8章。

\item[Embedding]
\textbf{嵌入/向量表示}。
将文本、图像等高维数据映射到低维连续向量空间的技术。
是向量数据库和 RAG 系统的基础。
参见第~\ref{ch:platform}~章向量数据库赛道分析。

\item[Fine-tuning]
\textbf{微调}。
在预训练模型的基础上，使用特定领域数据进行额外训练以适配下游任务。
包括全参数微调、LoRA、QLoRA 等方法。
参见第~\ref{ch:platform}~章微调工具赛道和 \textit{Emergence} 第8章。

\item[FLOPS]
\textbf{每秒浮点运算次数}（Floating-Point Operations Per Second）。
衡量计算设备算力的标准单位。AI 训练通常用 FP16/BF16 TFLOPS 衡量，
推理则关注 INT8/FP8 吞吐量。
参见第~\ref{ch:infrastructure}~章芯片性能对比。

\item[Foundation Model]
\textbf{基础模型}。
在大规模数据上预训练的大型 AI 模型，可通过微调或提示适配多种下游任务。
GPT-4、Claude、Llama 均属此类。
"基础模型"一词由斯坦福 HAI 于2021年提出。
参见第~\ref{ch:models}~章。

\item[GPU]
\textbf{图形处理器}（Graphics Processing Unit）。
最初为图形渲染设计的并行计算芯片，因其大规模并行架构
成为深度学习训练和推理的主力硬件。NVIDIA 控制90\%+市场。
参见第~\ref{ch:infrastructure}~章。

\item[Guardrails]
\textbf{护栏/安全围栏}。
限制 AI 模型输出不当内容（有害、虚假、偏见等）的技术和策略。
包括输入过滤、输出检测、内容策略执行等。
Lakera 等公司专注于此赛道。参见第~\ref{ch:platform}~章。

\item[Hallucination]
\textbf{幻觉}。
模型生成看似合理但实际错误或虚构的内容的现象。
是 LLM 应用落地的核心挑战之一，RAG 和 Grounding 是
主要缓解手段。参见第~\ref{ch:platform}~章和第~\ref{ch:patterns}~章。

\item[HBM]
\textbf{高带宽内存}（High Bandwidth Memory）。
采用垂直堆叠技术的高性能内存，为 GPU/AI 芯片提供远超传统 DRAM
的带宽。SK 海力士和三星是主要供应商。HBM 产能已成为 AI 芯片供应
链的关键瓶颈。参见第~\ref{ch:infrastructure}~章。

\item[Inference]
\textbf{推理}。
将训练好的模型部署到生产环境中处理新输入的过程。
推理的成本和延迟直接决定 AI 应用的商业可行性。
推理优化是基础设施层最活跃的创业赛道之一。
参见第~\ref{ch:infrastructure}~章§2.4。

\item[KV Cache]
\textbf{键值缓存}（Key-Value Cache）。
在自回归推理中缓存已计算的 Key 和 Value 矩阵，
避免重复计算。KV Cache 的内存占用随序列长度线性增长，
是长上下文推理的主要内存瓶颈。PagedAttention 等技术
旨在优化 KV Cache 管理。参见 \textit{Emergence} 第9章。

\item[LLM]
\textbf{大语言模型}（Large Language Model）。
基于 Transformer 架构、在大规模文本数据上预训练的语言模型，
参数规模通常在数十亿至数万亿之间。
GPT-4、Claude、Llama、千问等均属此类。
是本书讨论的核心技术对象。

\item[LoRA]
\textbf{低秩适配}（Low-Rank Adaptation）。
一种参数高效的微调方法，通过在预训练权重旁添加低秩矩阵来适配
新任务，只更新少量参数（通常< 1\%）。
大幅降低了微调的计算和存储成本。
参见 \textit{Emergence} 第8章。

\item[MoE]
\textbf{混合专家模型}（Mixture of Experts）。
一种模型架构，将参数分组为多个"专家"子网络，
每次推理只激活其中一部分。以更少的计算开销
实现更大的总参数量。Mixtral、DeepSeek-V3/R1、
Qwen-MoE 等均采用此架构。
参见 \textit{Emergence} 第6章。

\item[MLOps]
\textbf{机器学习运维}（Machine Learning Operations）。
将 DevOps 实践应用于机器学习系统的工程学科，
涵盖模型训练、部署、监控、版本管理的完整生命周期。
参见第~\ref{ch:platform}~章。

\item[Multimodal]
\textbf{多模态}。
能够同时处理和生成多种数据类型（文本、图像、音频、视频等）的
AI 系统。GPT-4o 和 Gemini 是多模态模型的代表。
从2024年开始成为基础模型竞赛的核心维度。

\item[Neocloud]
\textbf{新型 GPU 云}。
指2022年后兴起的一批以 GPU 算力为核心服务的云计算公司，
区别于传统三大云（AWS/Azure/GCP）。
CoreWeave、Lambda、Crusoe 等是代表。
参见第~\ref{ch:infrastructure}~章§2.1。

\item[NRR]
\textbf{净收入留存率}（Net Revenue Retention）。
SaaS 公司衡量现有客户收入变化的指标。
NRR > 100\% 表示现有客户的增购超过流失。
企业 AI 产品的 NRR 是判断 PMF 的关键信号。
参见第~\ref{ch:enterprise}~章。

\item[PagedAttention]
\textbf{分页注意力}。
vLLM 提出的 KV Cache 管理技术，借鉴操作系统虚拟内存的分页机制，
将 KV Cache 分割为固定大小的块（pages），实现非连续内存存储。
大幅减少内存碎片和浪费，是推理引擎的关键优化。
参见 \textit{Emergence} 第9章。

\item[PMF]
\textbf{产品市场匹配}（Product-Market Fit）。
产品满足市场需求的程度。在 AI 创业语境中，
PMF 的达成通常以稳定的收入增长、用户留存和 NRR 衡量。
参见第~\ref{ch:patterns}~章。

\item[Post-training]
\textbf{后训练}。
预训练完成后的所有优化步骤的统称，包括 SFT（监督微调）、
RLHF、DPO 等。后训练将"语言模型"转变为"AI 助手"。
参见 \textit{Emergence} 第8章。

\item[Pretraining]
\textbf{预训练}。
在大规模无标注数据上训练模型的初始阶段，
使模型习得语言理解和生成的基础能力。
预训练是 AI 产业链中资本密度最高的环节，
单次训练成本可达数亿美元。
参见 \textit{Emergence} 第5章。

\item[Prompt Engineering]
\textbf{提示工程}。
通过精心设计输入提示（Prompt）来引导模型产生期望输出的技术。
包括 few-shot、chain-of-thought、system prompt 等技巧。
是AI应用开发的基础技能。

\item[Quantization]
\textbf{量化}。
将模型权重和激活值从高精度（如 FP32/FP16）转换为低精度
（如 INT8/INT4）的技术。在可接受的精度损失范围内
显著减少模型的内存占用和计算量。
参见 \textit{Emergence} 第9章。

\item[RAG]
\textbf{检索增强生成}（Retrieval-Augmented Generation）。
将外部知识检索与模型生成相结合的技术范式。
先从知识库中检索相关文档，再将检索结果作为上下文输入模型，
有效缓解幻觉问题并提供可溯源的答案。
是平台层最核心的技术范式之一。
参见第~\ref{ch:platform}~章。

\item[RLHF]
\textbf{基于人类反馈的强化学习}（Reinforcement Learning from Human Feedback）。
通过人类评估员对模型输出的偏好排序来训练奖励模型，
再用强化学习优化生成策略。
ChatGPT 成功的关键技术之一。
参见 \textit{Emergence} 第8章。

\item[Scaling Law]
\textbf{缩放定律}。
描述模型性能与计算量、数据量、参数量之间幂律关系的经验规律。
Kaplan et al.\ (2020) 和 Chinchilla (Hoffmann et al., 2022)
是两个里程碑研究。缩放定律是"越大越好"范式的理论基础，
也是 AI 创业公司巨额融资的底层逻辑。
参见 \textit{Emergence} 第5章和本书第~\ref{ch:overview}~章。

\item[Speculative Decoding]
\textbf{推测解码}。
一种推理加速技术，使用小型"草稿"模型快速生成候选 token 序列，
再由大型模型并行验证。在不降低输出质量的前提下
提升推理吞吐量2--3倍。参见 \textit{Emergence} 第9章。

\item[SWE-bench]
\textbf{软件工程基准}（Software Engineering Benchmark）。
评估 AI 代码 Agent 自主解决真实 GitHub Issue 能力的基准测试。
从2024年的 $\sim$2\% 到2026年的 70\%+，
SWE-bench 分数的飙升是 AI 编码能力爆发的最佳度量。
参见第~\ref{ch:devtools}~章。

\item[Test-time Compute]
\textbf{推理时计算/测试时计算}。
在推理阶段投入额外计算以提升输出质量的技术范式。
o1、o3、DeepSeek-R1 等"推理模型"是其代表。
参见 \textit{Emergence} 第10章。

\item[Token]
\textbf{词元}。
模型处理文本的基本单位。一个 token 大约对应英文的 3/4 个单词
或中文的 1--2 个字。模型的输入输出均以 token 计量，
API 定价也以 token 为单位。
参见 \textit{Emergence} 第3章。

\item[TPU]
\textbf{张量处理单元}（Tensor Processing Unit）。
Google 自研的 AI 专用芯片。从 TPU v1（2016年）到 TPU v6e Trillium
（2025年），TPU 是最成功的非 NVIDIA AI 加速器。
参见第~\ref{ch:infrastructure}~章。

\item[Transformer]
\textbf{变换器}。
Vaswani et al.\ (2017) 提出的神经网络架构，
基于自注意力机制实现序列到序列的建模。
是 GPT、BERT、LLaMA 等几乎所有现代大模型的基础架构。
参见 \textit{Emergence} 第2章。

\item[Vector Database]
\textbf{向量数据库}。
专为存储和检索高维向量（Embedding）设计的数据库系统。
RAG 系统的核心组件。Pinecone、Weaviate、Qdrant、Milvus/Zilliz
是该赛道的主要玩家。
参见第~\ref{ch:platform}~章。

\item[vLLM]
一个高性能 LLM 推理引擎（开源），首创 PagedAttention 技术。
已成为 LLM 部署的事实标准推理框架之一。
参见 \textit{Emergence} 第9章和第~\ref{ch:infrastructure}~章。

\item[Wrapper]
\textbf{包装层/套壳应用}。
指在第三方模型 API 之上构建的应用，核心差异化主要来自 UX/UI
和 Prompt 设计，而非自有模型能力。"AI Wrapper"在2023--2024年
是最常见的 AI 创业形态，也是争议最大的类别——缺乏技术护城河
是其核心风险。
参见第~\ref{ch:patterns}~章。

\end{description}


% =============================================================
%  §A.3  DATA SOURCES AND METHODOLOGY
% =============================================================
\clearpage
\section{数据来源与方法论}
\label{sec:appendix-sources}

本书涵盖了2022年至2026年3月间全球 AI 创业生态的系统性调研。
以下说明数据来源、采集方法、验证流程和已知局限性。

\subsection{一手来源}

\begin{description}[style=nextline, leftmargin=1em]

\item[公司官方信息]
各公司官网、官方博客、技术文档、产品更新日志。
对于上市公司，参考其向 SEC（美国证券交易委员会）、
HKEX（香港交易所）、上交所/深交所提交的公开文件，
包括招股说明书、年报、季报。

\item[开源代码仓库]
GitHub、Hugging Face 上的开源项目仓库，包括代码提交历史、
Issue 讨论、Release Notes。特别关注 Star 数、Fork 数、
贡献者活跃度等量化指标。

\item[产品实测]
作者对本书涉及的主要产品（包括 Cursor、Claude Code、ChatGPT、
Perplexity、Midjourney、Gamma、Notion AI、各主要向量数据库等）
进行了实际使用和测试。产品评价基于截至2026年3月的版本。

\item[行业会议与访谈]
包括 NeurIPS 2024/2025、ICML 2025、AAAI 2026 等学术会议，
以及 Y Combinator Demo Day、TechCrunch Disrupt、极客公园
创新大会等行业活动中的公开演讲和 Panel 讨论。

\end{description}

\subsection{二手来源}

\begin{description}[style=nextline, leftmargin=1em]

\item[融资与估值数据]
Crunchbase、PitchBook、CB Insights 为全球融资数据的主要来源；
IT 桔子、烯牛数据为中国市场数据的主要来源。
对于估值和 ARR 数据，优先采用公司官方披露或 SEC Filing；
其次参考知名科技媒体（The Information、Bloomberg、Reuters）
的独家报道；最后采用数据库聚合数据。

\item[科技媒体]
英文来源：TechCrunch、The Information、Bloomberg Technology、
Ars Technica、The Verge、VentureBeat、Semafor。
中文来源：36氪、极客公园、量子位、机器之心、晚点 LatePost。
日韩来源：ITmedia、日経 xTECH、Byline Network。

\item[研究报告]
Bernstein Research、Morgan Stanley、Goldman Sachs 等投行研报；
McKinsey、BCG、Gartner、Forrester 等咨询机构报告；
斯坦福 HAI AI Index 年度报告、State of AI Report（Air Street Capital）、
EPOCH AI 等独立研究机构的数据与分析。

\item[社区与论坛]
Hacker News、Reddit（r/MachineLearning、r/LocalLLaMA）、
X/Twitter 上的 AI 社区讨论、Product Hunt 的产品发布数据。
中文社区包括知乎、即刻、V2EX。

\end{description}

\subsection{数据验证方法}

本书采用\textbf{多源交叉验证}原则：

\begin{enumerate}
  \item \textbf{三角验证}：关键数据点（融资额、估值、ARR、用户数等）
    至少需要两个独立来源相互印证。当来源之间存在差异时，
    优先采用离信息源最近的数据（如公司官方 > 投行研报 > 媒体报道），
    并在正文或脚注中注明差异。

  \item \textbf{时间戳标注}：所有数据均标注采集或引用的时间点。
    AI 创业领域数据变化极快，一家公司的估值可能在三个月内翻倍。
    因此，本书在涉及估值和融资数据时，尽可能注明具体的融资轮次日期。

  \item \textbf{利益冲突声明}：当引用来自投资方的分析报告时
    （如 a16z 对其被投公司的评价），会在上下文中注明潜在的
    利益关联，供读者自行判断。

  \item \textbf{公开可验证}：本书尽量采用公开可查的数据来源。
    对于来自匿名消息源或"知情人士"的信息，
    会在脚注中标注"据报道"或"据未经证实的消息"。
\end{enumerate}

\subsection{时效性说明}

\begin{itemize}
  \item \textbf{数据截止日期}：2026年3月21日。
    在此日期之后发生的融资、产品发布、上市等事件未纳入本书。

  \item \textbf{快速变化的领域}：AI 创业生态变化极为迅速。
    本书出版时，部分公司可能已经完成新一轮融资、估值显著变化、
    或甚至已经关闭。读者在引用具体数据时，建议查阅最新来源进行
    核实。

  \item \textbf{版本更新}：本书电子版将不定期更新关键数据。
    每次更新将在变更日志中记录修改内容和时间。
\end{itemize}

\subsection{已知局限性}

\begin{enumerate}
  \item \textbf{估值数据的不透明性}：大量 AI 创业公司为私营企业，
    其估值和收入数据来自融资轮次的媒体报道，而非公开审计。
    实际财务状况可能与报道数字存在偏差。

  \item \textbf{区域覆盖不均}：硅谷和中国的覆盖深度显著高于其他地区。
    中东、东南亚、拉美等新兴 AI 生态的覆盖可能存在遗漏。

  \item \textbf{语言偏差}：英语和中文语种的信息来源占绝对主导；
    日语、韩语、希伯来语、阿拉伯语来源的覆盖有限，
    可能低估了非英中语种生态的规模和活力。

  \item \textbf{幸存者偏差}：本书倾向于关注成功融资或获得媒体曝光的
    公司。大量未获报道但可能同样值得关注的早期创业公司未被收录。

  \item \textbf{对公领域的局限}：涉及国防、情报等领域的 AI 应用
    因保密性质，在本书中几乎未涉及。

  \item \textbf{作者立场}：作者的学术背景和地理位置（香港科技大学）
    可能影响对某些话题的视角和判断。读者宜结合多种来源形成
    自己的独立判断。
\end{enumerate}


% =============================================================
%  §A.4  CROSS-REFERENCE TO EMERGENCE
% =============================================================
\clearpage
\section{与 \textit{Emergence} 的交叉索引}
\label{sec:appendix-emergence-xref}

本书（\textit{The AI Startup Landscape}，以下简称"本卷"）
与姊妹卷 \textit{Emergence: The Technical Foundations of Large Language Models}
（以下简称"\textit{Emergence}"）共同构成 LLM Observatory Monograph Series
的两卷本体系。\textit{Emergence} 聚焦技术原理（"How it works"），
本卷聚焦商业生态（"Who builds on it and how they make money"）。

以下表格列出两卷之间的章节对应关系和关联主题，
方便读者在技术与商业之间快速跳转阅读。

\bigskip

\begin{tabularx}{\textwidth}{@{}l l X@{}}
\toprule
\rowcolor{figLightGray}
\textbf{本卷章节} & \textbf{\textit{Emergence} 章节} & \textbf{关联主题} \\
\midrule

Ch.\,\ref{ch:overview} 全景概述 &
Ch.\,1 为什么"预测下一个词"能产生智能？ &
技术驱动力概述：Scaling Laws 如何催生创业浪潮 \\
\addlinespace

Ch.\,\ref{ch:infrastructure} 基础设施层 &
Ch.\,9 推理优化 &
GPU 云/推理芯片的商业逻辑 $\leftrightarrow$ 推理瓶颈的技术原理 \\
\addlinespace

Ch.\,\ref{ch:infrastructure} 基础设施层 &
Ch.\,5 预训练 &
训练集群商业模式 $\leftrightarrow$ 预训练的计算需求和工程挑战 \\
\addlinespace

Ch.\,\ref{ch:infrastructure} 基础设施层 &
Ch.\,6 MoE &
芯片架构需求 $\leftrightarrow$ MoE 模型对通信带宽和内存的特殊需求 \\
\addlinespace

Ch.\,\ref{ch:models} 模型层 &
Ch.\,2 Transformer &
模型公司的技术路线选择 $\leftrightarrow$ Transformer 架构演进 \\
\addlinespace

Ch.\,\ref{ch:models} 模型层 &
Ch.\,5 预训练 &
预训练成本的商业影响 $\leftrightarrow$ 数据配比和训练策略 \\
\addlinespace

Ch.\,\ref{ch:models} 模型层 &
Ch.\,8 后训练 &
RLHF/DPO 对产品差异化的影响 $\leftrightarrow$ 对齐技术细节 \\
\addlinespace

Ch.\,\ref{ch:models} 模型层 &
Ch.\,10 推理时计算 &
推理模型（o1/o3/R1）的商业价值 $\leftrightarrow$ Test-time Compute 理论 \\
\addlinespace

Ch.\,\ref{ch:platform} 平台层 &
Ch.\,7 长上下文 &
RAG vs.\ 长上下文的产品选择 $\leftrightarrow$ 长上下文技术实现 \\
\addlinespace

Ch.\,\ref{ch:platform} 平台层 &
Ch.\,9 推理优化 &
推理引擎商业化（vLLM/TensorRT-LLM） $\leftrightarrow$ PagedAttention 等技术 \\
\addlinespace

Ch.\,\ref{ch:devtools} 开发者工具 &
Ch.\,11 评估与安全 &
SWE-bench 商业影响 $\leftrightarrow$ 代码评估方法论 \\
\addlinespace

Ch.\,\ref{ch:devtools} 开发者工具 &
Ch.\,15 Agent 大战编年史 &
代码 Agent 产品竞赛 $\leftrightarrow$ Agent 架构和能力演进 \\
\addlinespace

Ch.\,\ref{ch:consumer-create} 创意与内容生成 &
Ch.\,2 Transformer &
扩散模型商业化 $\leftrightarrow$ 多模态 Transformer 架构 \\
\addlinespace

Ch.\,\ref{ch:consumer-social} AI 伴侣与搜索 &
Ch.\,7 长上下文 &
AI 伴侣的"记忆"功能 $\leftrightarrow$ 长上下文技术支撑 \\
\addlinespace

Ch.\,\ref{ch:consumer-social} AI 伴侣与搜索 &
Ch.\,8 后训练 &
对话体验差异化 $\leftrightarrow$ 人格化对齐技术 \\
\addlinespace

Ch.\,\ref{ch:enterprise} 企业 AI &
Ch.\,8 后训练 &
垂直领域微调策略 $\leftrightarrow$ SFT/RLHF 工程实践 \\
\addlinespace

Ch.\,\ref{ch:china} 中国 AI 生态 &
Ch.\,12 团队与组织 &
中国 AI 人才管线 $\leftrightarrow$ 全球 AI 研究团队分布 \\
\addlinespace

Ch.\,\ref{ch:china} 中国 AI 生态 &
Ch.\,6 MoE &
DeepSeek 技术路线 $\leftrightarrow$ MoE 架构创新细节 \\
\addlinespace

Ch.\,\ref{ch:global} 全球 AI 生态 &
Ch.\,12 团队与组织 &
各国 AI 生态比较 $\leftrightarrow$ 全球研究实验室格局 \\
\addlinespace

Ch.\,\ref{ch:funding} 资本与融资 &
Ch.\,5 预训练 &
巨额融资的底层逻辑 $\leftrightarrow$ 预训练成本的数量级分析 \\
\addlinespace

Ch.\,\ref{ch:patterns} 规律与教训 &
Ch.\,13 前沿展望 &
商业模式演进规律 $\leftrightarrow$ 技术趋势对商业的影响 \\
\addlinespace

Ch.\,\ref{ch:outlook} 展望 &
Ch.\,13 前沿展望 &
2026下半年商业预判 $\leftrightarrow$ 技术路线图前瞻 \\
\addlinespace

Ch.\,\ref{ch:outlook} 展望 &
Ch.\,15 Agent 大战编年史 &
Agent 经济的商业机会 $\leftrightarrow$ Agent 技术演进时间线 \\

\bottomrule
\end{tabularx}

\bigskip

\noindent\textbf{阅读建议：}

\begin{itemize}
  \item \textbf{技术优先路径}：先读 \textit{Emergence} 全书以建立技术理解，
    再读本卷以了解这些技术如何在商业世界中落地。适合工程师和研究者。

  \item \textbf{商业优先路径}：先读本卷以了解 AI 创业的全景，
    遇到技术概念时参考术语表（§\ref{sec:appendix-glossary}）
    或跳转至 \textit{Emergence} 对应章节。适合投资人和产品经理。

  \item \textbf{主题优先路径}：根据个人兴趣，利用上表选择一个主题
    （如"推理优化"），同时阅读两卷对应章节，建立技术-商业的
    完整认知。适合聚焦特定赛道的从业者。
\end{itemize}
